\documentclass[12pt]{article}

\usepackage[utf8]{inputenc}
\usepackage[T2A]{fontenc}
\usepackage[russian]{babel}
\usepackage{url}

\usepackage{amsmath, amssymb}
\usepackage{physics}
\usepackage{geometry}
\usepackage{graphicx}
\usepackage{setspace}

\geometry{margin=2.5cm}
\onehalfspacing

\title{Physics-Informed Neural Networks (PINN)}
\author{}
\date{}

\begin{document}

\section{Physics-Informed Neural Networks для задач термоупругой волновой динамики}

\subsection{Общая идея метода PINN}

Physics-Informed Neural Networks (PINN) представляют собой класс нейросетевых моделей, в которых априорные физические законы вводятся непосредственно в процесс обучения через функцию потерь. 

В отличие от классических методов машинного обучения, ориентированных исключительно на аппроксимацию данных, PINN позволяют решать задачи математической физики, формулируемые в виде дифференциальных уравнений в частных производных.

Основная идея метода заключается в аппроксимации искомых физических полей нейросетью
\begin{equation}
\mathcal{N}_\theta : (x,t) \mapsto \bigl( u_i(x,t), T(x,t) \bigr),
\end{equation}
где $u_i$ — компоненты вектора смещений, $T$ — температурное поле, а $\theta$ — параметры нейросети.

Производные по пространству и времени вычисляются с использованием автоматического дифференцирования, что позволяет включать остатки дифференциальных уравнений непосредственно в функционал потерь.

\subsection{Применение PINN к системе связанных уравнений термоупругости}

Рассматривается система линейной термоупругости Куна–Томсона, включающая уравнение движения
\begin{equation}
\rho \frac{\partial^2 u_i}{\partial t^2} = \frac{\partial \sigma_{ij}}{\partial x_j},
\end{equation}
определяющее динамику упругих волн, и уравнение теплопереноса
\begin{equation}
k \nabla^2 T - C \frac{\partial T}{\partial t} = \gamma T_0 \frac{\partial \varepsilon_{kk}}{\partial t},
\end{equation}
учитывающее термомеханическую генерацию тепла.

Связь между напряжениями и деформациями задаётся соотношением
\begin{equation}
\sigma_{ij} = \lambda \delta_{ij} \varepsilon_{kk} + 2\mu \varepsilon_{ij} - \gamma \delta_{ij} T,
\end{equation}
где $\lambda$, $\mu$ — параметры Ламе, $\varepsilon_{ij}$ — тензор деформаций.

В рамках PINN нейросеть аппроксимирует одновременно все искомые поля $u_i(x,t)$ и $T(x,t)$, а физические уравнения используются в качестве ограничений.

\subsection{Формирование функции потерь}

Функция потерь PINN строится как взвешенная сумма невязок дифференциальных уравнений и начально-граничных условий:
\begin{equation}
\mathcal{L} =
\mathcal{L}_{\text{motion}}
+ \mathcal{L}_{\text{heat}}
+ \mathcal{L}_{\text{BC}}
+ \mathcal{L}_{\text{IC}}.
\end{equation}

Невязка уравнения движения определяется как
\begin{equation}
\mathcal{L}_{\text{motion}} =
\left\|
\rho \frac{\partial^2 u_i}{\partial t^2}
- \frac{\partial \sigma_{ij}}{\partial x_j}
\right\|^2,
\end{equation}
а невязка уравнения теплопереноса — как
\begin{equation}
\mathcal{L}_{\text{heat}} =
\left\|
k \nabla^2 T
- C \frac{\partial T}{\partial t}
- \gamma T_0 \frac{\partial \varepsilon_{kk}}{\partial t}
\right\|^2.
\end{equation}

Начальные и граничные условия вводятся в виде дополнительных штрафных слагаемых, что обеспечивает физическую корректность решения.

\subsection{Учет неоднородной геологической среды}

Пространственная неоднородность геологической среды учитывается за счёт зависимости параметров $\rho(x)$, $\lambda(x)$, $\mu(x)$, $k(x)$ и $\gamma(x)$ от координат. Эти параметры могут быть заданы в табличной форме и интерполированы либо непосредственно поданы на вход нейросети:
\begin{equation}
\mathcal{N}_\theta : (x,t,\rho,\lambda,\mu,k,\gamma) \mapsto (u_i,T).
\end{equation}

Таким образом, PINN естественным образом обрабатывает вариации физических характеристик среды без необходимости перестроения вычислительной сетки.

\subsection{Алгоритм решения задачи с использованием PINN}

Процедура решения задачи термоупругой волновой динамики с применением PINN включает следующие этапы:
\begin{enumerate}
    \item Формирование обучающей выборки коллокационных точек в пространственно-временной области.
    \item Задание нейросетевой аппроксимации искомых полей.
    \item Автоматическое вычисление пространственно-временных производных.
    \item Формирование функции потерь на основе физических уравнений и граничных условий.
    \item Оптимизация параметров нейросети методом градиентного спуска.
\end{enumerate}

В результате обучения нейросеть предоставляет непрерывное приближение решения системы связанных уравнений термоупругости во всей области определения.

\subsection{Преимущества метода}

Применение PINN для задач термоупругих волн в геологических средах обладает рядом преимуществ:
\begin{itemize}
    \item строгий учет физических законов;
    \item возможность работы с малым объемом экспериментальных данных;
    \item отсутствие необходимости построения пространственной сетки;
    \item удобство моделирования неоднородных и анизотропных сред.
\end{itemize}


\begin{thebibliography}{9}

\bibitem{raissi2019}
M.~Raissi, P.~Perdikaris, G.~E.~Karniadakis,
\textit{Physics-informed neural networks: A deep learning framework for solving forward and inverse problems involving nonlinear partial differential equations},
Journal of Computational Physics, Vol.~378, 2019, pp.~686--707.\\
DOI: \url{https://doi.org/10.1016/j.jcp.2018.10.045}\\
arXiv: \url{https://arxiv.org/abs/1711.10561}

\bibitem{karniadakis2021}
G.~E.~Karniadakis, I.~G.~Kevrekidis, L.~Lu, P.~Perdikaris, S.~Wang, L.~Yang,
\textit{Physics-informed machine learning},
Nature Reviews Physics, Vol.~3, 2021, pp.~422--440.\\
DOI: \url{https://doi.org/10.1038/s42254-021-00314-5}\\
arXiv: \url{https://arxiv.org/abs/2104.13420}

\bibitem{bui2021}
T.~Q.~Bui, C.~Zhang,
\textit{Physics-informed neural networks for computational mechanics},
Engineering Analysis with Boundary Elements, Vol.~129, 2021, pp.~160--175.\\
DOI: \url{https://doi.org/10.1016/j.enganabound.2021.03.006}

\bibitem{nowacki1986}
W.~Nowacki,
\textit{Thermoelasticity},
Pergamon Press, Oxford, 1986.\\
URL: \url{https://www.sciencedirect.com/book/9780080105919/thermoelasticity}

\bibitem{auld1990}
B.~A.~Auld,
\textit{Acoustic Fields and Waves in Solids},
Wiley, New York, 1990.\\
URL: \url{https://onlinelibrary.wiley.com/doi/book/10.1002/9780470172834}

\end{thebibliography}

\end{document}